\section{Introduction}

The SO-101 robot arm is commonly operated through a host-computer-based software stack, in which motion commands are transmitted to the robot through an external communication interface \cite{lerobot_so101}. While this approach simplifies development, it also creates a permanent dependency on a computer for communication, kinematic calculations, and high-level control. Executing these functions directly on a microcontroller can reduce this dependency and provides a foundation for compact, standalone robotic systems. However, embedded implementation requires reliable servo communication, computationally efficient kinematics, safe handling of joint limits, and integration with external sensors.

This work presents \textit{elroboto}, a proof-of-concept embedded control system for the SO-101 robot arm. An STM32U545 microcontroller communicates directly with the daisy-chained Feetech servos and executes the main motion-control functions without a host computer. Four joints are used as an active kinematic chain for Cartesian position control, while wrist rotation and gripper actuation are treated as fixed degrees of freedom. A forward-kinematic model derived from the robot description is combined with a numerical damped least-squares inverse-kinematic solver to convert Cartesian target positions into valid servo commands.

The developed system additionally includes servo feedback, communication-error handling, calibrated joint limits, and blocking and non-blocking motion interfaces. The computational cost of the kinematic implementation is evaluated by benchmarking alternative trigonometric functions and homogeneous-transformation multiplication methods. The motion stack is validated on the physical robot through a square-shaped Cartesian waypoint sequence. In parallel, a wireless distance-sensing subsystem is developed as the basis for collision-feedback integration. The presented system is intended as a proof of concept and does not provide full pose control, general path planning, or guaranteed Cartesian trajectory tracking.